\documentclass[a4paper, 12pt]{article}

\usepackage[T1]{fontenc}      % Supports accented characters
\usepackage[french]{babel}
\usepackage{textcomp}
\usepackage{amssymb}
 \usepackage{newtxtext} \usepackage{newtxmath}
\usepackage{xcolor}
\usepackage{tcolorbox}
\usepackage{expex}
\usepackage{forest}
\DeclareMathAlphabet{\mathcal}{OMS}{cmsy}{m}{n}
\usepackage{parskip}
\newcommand{\verbf}[1]{\textcolor{blue}{#1}}
\newcommand{\subjf}[1]{\underline{#1}}
\newcommand{\objf}[1]{\textcolor{red}{#1}}
\newcommand{\tensef}[1]{\colorbox{yellow}{#1}}
\newcommand{\noun}[1]{\textsc{#1}}
\newcommand{\adj}[1]{\textit{#1}}
\newcommand{\prep}[1]{\textcolor{purple}{#1}}
\usepackage{graphicx}
\begin{document}


\begin{itemize}
    \item Mi (masculino) $\to $ Mon 
    \item Mi (femenino) $\to $ ma
    \item Mis (plural) $\to $ mes
    \item Son (masculino) $\to $ su 
    \item Sa (femenino) $\to $ su 
    \item Ses (plural) $\to $ sus
\end{itemize}

\begin{table}[htpb]
\centering
\begin{tabular}{|l|l|l|l|}
\hline
\textbf{Pronombre} & \textbf{Présent} & \textbf{Imparfait} & \textbf{Futur simple} \\ \hline
Je / J' & suis & étais & serai \\ \hline
Tu & es & étais & seras \\ \hline
Il / Elle / On & est & était & sera \\ \hline
Nous & sommes & étions & serons \\ \hline
Vous & êtes & étiez & serez \\ \hline
Ils / Elles & sont & étaient & seront \\ \hline
\end{tabular}
\caption{Conjugación del verbo \textit{être} (Tiempos simples del Indicativo)}
\end{table}

\vspace{0.5cm}

\begin{table}[htpb]
\centering
\begin{tabular}{|l|l|l|l|}
\hline
\textbf{Pronombre} & \textbf{Passé composé} & \textbf{Conditionnel prés.} & \textbf{Subjonctif prés.} \\ \hline
J' & ai été & serais & sois \\ \hline
Tu & as été & serais & sois \\ \hline
Il / Elle / On & a été & serait & soit \\ \hline
Nous & avons été & serions & soyons \\ \hline
Vous & avez été & seriez & soyez \\ \hline
Ils / Elles & ont été & seraient & soient \\ \hline
\end{tabular}
\caption{Conjugación del verbo \textit{être} (Pasado compuesto, Condicional y Subjuntivo)}
\end{table}

\vspace{0.5cm}

\begin{table}[htpb]
\centering
\begin{tabular}{|l|l|}
\hline
\textbf{Pronombre (Implícito)} & \textbf{Impératif présent} \\ \hline
(Tu) & sois \\ \hline
(Nous) & soyons \\ \hline
(Vous) & soyez \\ \hline
\end{tabular}
\caption{Conjugación del verbo \textit{être} (Imperativo)}
\end{table}

\begin{table}[htpb]
\centering
\begin{tabular}{|l|l|l|l|}
\hline
\textbf{Pronombre} & \textbf{Passé composé} & \textbf{Conditionnel prés.} & \textbf{Subjonctif prés.} \\ \hline
J' / Je* & ai eu & aurais & aie \\ \hline
Tu & as eu & aurais & aies \\ \hline
Il / Elle / On & a eu & aurait & ait \\ \hline
Nous & avons eu & aurions & ayons \\ \hline
Vous & avez eu & auriez & ayez \\ \hline
Ils / Elles & ont eu & auraient & aient \\ \hline
\end{tabular}
\caption{Conjugación del verbo \textit{avoir} (Pasado compuesto, Condicional y Subjuntivo)}
\end{table}

\begin{table}[htbp]
\centering
\caption{Conjugación de verbos regulares del primer grupo (-er) | Modelo: \textit{parler}}
\vspace{0.2cm} % Espacio entre el título y la tabla
\begin{tabular}{lcccccc}
\toprule
\textbf{Tiempo / Modo} & \textbf{je (j')} & \textbf{tu} & \textbf{il/elle/on} & \textbf{nous} & \textbf{vous} & \textbf{ils/elles} \\
\midrule
\textbf{Présent} & parle & parles & parle & parlons & parlez & parlent \\
\textbf{Imparfait} & parlais & parlais & parlait & parlions & parliez & parlaient \\
\textbf{Passé simple} & parlai & parlas & parla & parlâmes & parlâtes & parlèrent \\
\textbf{Futur simple} & parlerai & parleras & parlera & parlerons & parlerez & parleront \\
\textbf{Conditionnel} & parlerais & parlerais & parlerait & parlerions & parleriez & parleraient \\
\textbf{Passé composé} & ai parlé & as parlé & a parlé & avons parlé & avez parlé & ont parlé \\
\bottomrule
\end{tabular}
\end{table}
\begin{table}[htbp]
\centering
\caption{Conjugación de verbos regulares del segundo grupo (-ir) | Modelo: \textit{finir}}
\vspace{0.2cm}
\begin{tabular}{lcccccc}
\toprule
\textbf{Tiempo / Modo} & \textbf{je (j')} & \textbf{tu} & \textbf{il/elle/on} & \textbf{nous} & \textbf{vous} & \textbf{ils/elles} \\
\midrule
\textbf{Présent} & finis & finis & finit & finissons & finissez & finissent \\
\textbf{Imparfait} & finissais & finissais & finissait & finissions & finissiez & finissaient \\
\textbf{Passé simple} & finis & finis & finit & finîmes & finîtes & finirent \\
\textbf{Futur simple} & finirai & finiras & finira & finirons & finirez & finiront \\
\textbf{Conditionnel} & finirais & finirais & finirait & finirions & finiriez & finiraient \\
\textbf{Passé composé} & ai fini & as fini & a fini & avons fini & avez fini & ont fini \\
\bottomrule
\end{tabular}
\end{table}

\begin{table}[htbp]
\centering
\caption{Conjugación de verbos regulares en -oir | Modelo aproximado: \textit{recevoir}}
\vspace{0.2cm}
\begin{tabular}{lcccccc}
\toprule
\textbf{Tiempo / Modo} & \textbf{je (j')} & \textbf{tu} & \textbf{il/elle/on} & \textbf{nous} & \textbf{vous} & \textbf{ils/elles} \\
\midrule
\textbf{Présent} & reçois & reçois & reçoit & recevons & recevez & reçoivent \\
\textbf{Imparfait} & recevais & recevais & recevait & recevions & receviez & recevaient \\
\textbf{Passé simple} & reçus & reçus & reçut & reçûmes & reçûtes & reçurent \\
\textbf{Futur simple} & recevrai & recevras & recevra & recevrons & recevrez & recevront \\
\textbf{Conditionnel} & recevrais & recevrais & recevrait & recevrions & recevriez & recevraient \\
\textbf{Passé composé} & ai reçu & as reçu & a reçu & avons reçu & avez reçu & ont reçu \\
\bottomrule
\end{tabular}
\end{table}

\pagebreak

$\S$ Rimbaud, dans la cinquième partie du poème \textit{Enfance}, s'identifie à
quatre choses : un saint, un savant, un piéton, un enfant abandonné qui prend la
mer. Je suis aussi toutes ces choses. Je suis aussi parti. J'ai longtemps vu le
visage d'une femme qui n'existe plus. J'ai été aussi un saint sous l'ermitage de
ses cheveux. Je suis aussi un savant, car j'ai lu dans ses yeux le livre de la
vie. Je suis un piéton sur le chemin d'un rêve rêvé par elle. Et je suis la
longue rumeur de la rivière où nous avons plongé.

Paulina, tandis que le monde t'oublie, je me souviens de toi. Je garde ton
visage, le souvenir de la rivière où nous étions enfants, où nous étions
heureux, où l'eau t'enveloppait. Je t'ai dit : «Comme l'eau, telle sera ma vie,
telle sera ta vie, et un jour nous n'existerons plus». Et toi, tu m'as
éclaboussé, pour que je cesse d'être si sot.
~

$\S$ Aujourd'hui, je vais étudier les langages formels appelés
$\omega$-réguliers. Ces langages sont les langages acceptés par les automates de
Büchi, qui sont une chose très belle. Je me suis fait un café. C'est nuageux. Il
fait froid. Je suis heureux.

~ 

$\S$ A few months ago I wrote a brief commentary on Nietzsche's work \textit{On
truth and lies in a non-moral sense}. An example of the seemingly infinite
expressive power of poetry is the fact that the total content of that essay, and
of my comments on it, is perfectly summarized in the opening verses of
Mallarmé's \textit{Salut}:

\begin{quote}
    Rien, cette écume, vierge verse \newline 
    A né désigner que la coupe.
\end{quote}

~ 

$\S$ Tandis que le monde t'oublie, je me souviens de toi. Je garde ton visage, le
souvenir de la rivière où nous étions enfants, où nous étions heureux, où l'eau
t'enveloppait. Je t'ai dit : «Comme l'eau, telle sera ma vie, telle sera ta
vie, et un jour nous n'existerons plus». Et toi, tu m'as éclaboussé, pour que je
cesse d'être si sot.

~ 



$\S$ \textit{Cauchemar.} 
J'ai fait un cauchemar très inquiétant. Mon père devait
m'administrer l'injection létale. Tandis que j'étais allongé sur le lit où,
autrefois, lui et ma mère dormaient, dans l'ancienne maison familiale, il m'a
administré la mort. Je sentais la mort comme un assoupissement graduel qui
commençait dans mes bras et mes jambes, une peur subtile qui était à la fois
physique et spirituelle. Au début, je sentais une résignation pacifique, mais quand
la pensée du S s'est formée dans mon cœur, le désespoir a pris possession de moi.
J'ai couru vers la rue, la mort fleurissant encore en moi, dans l'espoir de la voir,
de lui dire que bien que ma vie se flétrissait, elle avait fleuri pour elle un jour.
Et que ce jour-là, j'ai été heureux.

~ 

$\S$ De femme á femme, qu'est-ce que je cherche? Je me sens comme un naufragé
dans une mer de solitude. Je suis mort dans un baiser cette nuit, mais je sens
que rien ne peut étre perfait, que rien n'est beau, que je vis une vie de
vacuité totale. Le néant est la couleur du baiser que je cherche: je désire le
néant, bien que je l'ignore. On donne beaucoup de baisers dans la vie, mais un
seul ou deux sont capables d'assassiner la nostalgie absolue de la vie.

~ 

$\S$ Aux heures de tristesse, je suis comme une fleur qui lentement couvre le
monde de beauté. La tristesse, me semble-t-il, est la facilitatrice de l'espace esthétique. Des
égouts, elle fait ses pétales. Peut-être est-elle la main qui découvre le voile
de fadeur qui cache et dissimule notre singularité. C'est pourquoi, demoiselle
bleue, je te cherche et te perçois en toutes les choses. Je t'étreins et
t'enlace, je t'embrasse tendrement, ma tristesse, nuit sans fin!

~ 

$\S$ Les souvenirs sont des lucioles qui s'allument et s'éteignent dans l'obscurité
de la conscience. Et, lentement, nous avons semé un champ de lucioles. Et
maintenant qu'elles dansent et se mêlent, tandis que la lumière de l'une ou
de l'autre peut paraître languissante ou mélancolique, l'ensemble est
lumineux et incandescent. En les regardant, je sens que cette histoire a atteint
un point d'équilibre. Nous orbitons, chacun à la même distance, autour d'un
centre secret. Personne ne le sait, mais c'est ainsi. Je pense: «Je partirai
bientôt pour un autre pays, et l'amour que je lui ai offert s'éteindra». Mais c'est
un mensonge. Tout comme le Paraná poursuivra son dénouement aveugle vers le cœur de
l'homme, sur les eaux obscures, quand la lune étonnera l’œil des poissons,
on verra encore le reflet frémissant des lucioles.

~ 

$\S$ Tu ne vieillis plus. Ton visage, dans ma mémoire, est le même visage que
j'ai embrassé un jour, la même fleur fanée de ta jeunesse. Et, bien que je
vieillisse, sache que je marche encore avec toi. Je veux être nu avec toi, et je ne
parle pas de mon corps, et je ne parle pas de tes ossements. Combien d'années se
sont écoulées? J'ai flotté, j'ai été entraîné toutes ces années. Tu sais?
J'imagine encore, quand je marche dans une rue quelconque, que tu marches avec
moi. Je crois entendre ta voix encore. Chaque nuit, je m'endors en te ramenant à
la vie.

~ 
\pagebreak

$\S$ Quand j'avais dix-sept ans, je visitai le temple jaïniste de Ranakpur. J'y
appris l'existence d'une pratique encore répandue aujourd'hui. Les moines
jaïnistes mènent une vie d'ascétisme silencieux. Certains d'entre eux, ayant
atteint l'émancipation suprême, s'abandonnent à la méditation et à la
désintégration par inanition.

Durant le même voyage, je visitai le fort de Jaisalmer, érigé à l'aube du
deuxième millénaire après Jésus-Christ. Au XIIIe siècle, tandis que les femmes
commettaient le jauhar, les quelques hommes survivants, succombant au siège de
l'envahisseur Alâ ud-Dîn, quittèrent le fort pour mourir au combat.

Gödel, qui était paranoïaque, mourut d'inanition par peur d'être empoisonné.
Mishima, cet écrivain exquis, commit le seppuku lorsque son coup d'État échoua.
Caton d'Utique rechercha la mort d'une manière semblable : vaincu par Jules
César, il plongea son épée dans sa poitrine et déchira ses entrailles avec ses
propres mains. Sénèque mourut d'une mort lente et bureaucratique sur ordre de
Néron. L'agent bolivien qui assassina le Che Guevara raconta que ses derniers
mots furent: «Calmez-vous et visez bien: vous allez tuer un homme». Severino Di
Giovanni, faisant face à l'exécution, dit à son avocat: «Je suis conscient de
ma situation et je n'ai pas l'intention de me dérober à quelque responsabilité
que ce soit. J'ai joué, j'ai perdu. Comme tout bon perdant, je paie le prix de
ma vie». Socrate, par respect pour un verdict qu'il savait injuste, refusa
d'être sauvé par ses amis. Le moment venu, il but la ciguë calmement, et sa
dernière pensée fut à propos d'une dette qu'il ne voulait pas laisser impayée.

Au cœur de l'\textit{ethos} chrétien se trouve la volonté de mourir. Dieu
lui-même s'incarna pour souffrir une mort terrestre. Origène souffrit deux annés
de torture aux mains de Dèce: il jugea la mort un moindre mal que de renier sa
foi. La Bible est du même avis: Dans l'Évangile selon Marc, Pierre dit à Jésus:
«Quand il me faudrait mourir avec toi, je ne te renierai pas». Dans les
\textit{Samuel} 31:4, on nous dit que Saül, vaincu par les Philistins, dit à son
écuyer:

\begin{quote}
Tire ton épée, et m'en transperce, de peur que ces incirconcis ne viennent me
percer et me faire subir leurs outrages. Celui qui portait ses armes ne voulut
pas, car il était saisi de crainte. Et Saül prit son épée, et se jeta dessus.
\end{quote}

Hilaire de Poitiers, quand il était en Syrie, sut que sa fille était courtisée
par des hommes très riches et puissants. Inspiré par l'idée qu'aucun mariage
n'était préférable pour sa fille que le mariage avec Dieu, il pria pieusement
pour sa mort. Curieusement, elle mourut, et il se consola de sa mort. Son
épouse, qui trouva aussi du bonheur en la mort de sa fille, lui demanda de prier
pour sa propre mort. Et, peu après, elle mourut d'une mort qui les satisfit tous
les deux.

Bien que nous sentions une certaine inclination à considérer la mort volontaire
comme un mal, c'est en fait surprenant avec quelle résolution une personne peut
renoncer à la volonté de vivre ou se consacrer à la volonté de mourir. Quand
j'étais jeune, me surprenait la facilité avec laquelle certains organismes
peuvent renoncer à la volonté de perpétuer leurs propres vies. (Je ne comprenais
pas la théorie de l'évolution.) Il est fascinant que la nature soit assez
complexe pour produire des agents qui semblent la contredire. Sur ce point, elle
ressemble à un créateur qui confère à sa création la faculté de se rebeller,
concept qui est la racine du mythe chrétien, comme le dit Milton:

\begin{quote}
Because we freely love, as in our will\newline
To love or not; in this we stand or fall:\newline
And some are fall'n, to disobedience fall'n (...).
\end{quote}

Bien des destinées sont pires que la mort, et la mort n'est pas le pire des
malheurs, si toutefois elle est un malheur. Lucrèce soutenait, de manière
irréfutable, que nos sentiments par rapport à la mort devraient être similaires
à ceux que nous éprouvons à l'égard de notre inexistence avant de naître. Lucain
considérait la mort comme si précieuse qu'il estimait injuste que tout le monde
puisse mourir. L'immolation collective est tragique et grotesque, mais il paraît
indiscutable que le destin qui attendait les femmes et les enfants du fort de
Jaisalmer était encore pire. C'est vrai que la plupart d'entre nous éprouvent un
attachement à la vie, et que nous ne désirons pas «go gentle into that good
night». Mais les exemples que j'ai donnés montrent qu'il y a place pour la paix,
y compris la satisfaction, dans ce crépuscule éternel.




\begin{itemize}
    \item «Llevan adelante», «to conduct, to carry out» $\to $ mènent.
    \item \textit{Lorsque} es más formal que \textit{quand}, pero ambos son
        correctos para expresar que una cosa pasó \textit{cuando} otra pasó. 
        Pero \textit{lorsque} no se puede usar en una pregunta, mientras que 
        \textit{quand} sí. 
    \item Sa épouse $\to $ su esposa con esposa \textit{nf}. Pero por eufonía se
        dice \textit{son épouse} como si fuera si fuera \textit{nm}.
    \item Trouver du bonheur: hallar cierta felicidad. Trouver la bonheur:
        hallar la felicidad. Pero bonheur necesita un artículo o un partitivo.
    \item Después de aucun, el francés vuelve a negar el verbo. «Ningún
        matrimonio era ...» $\to $ «Aucun mariage n'était ...» $\approx$ «Ningún
        matrimonio no era ...». 
    \item «Es <adj> con cuánto/a <sustantivo> <suj> <ver>» se expresa sin usar
        «cuánto» (combien) sino quelle (cuál, qué). E.g. «Es extraño con
        cuánta locura (él) vive» $\to $ «C'est étrange avec quelle folie il
        vit».
\end{itemize}


\pagebreak 

$\S$ 
Si j'ai vu le cercueil, la pierre froide, l'herbe fleurie, la mère en
pleurs; enfin, les signes universels de ta réalité finale. Si je me suis allongé
sur la rosée qui tapisse ta tombe, sur un lit qui ne fait qu'un avec tes
cauchemars, sur un désir sans objet ni sens. Si j'ai hésité, comme Jésus-Christ;
si j'ai consacré ma chair à des tigres affamés, comme Siddhartha; si enfin je ne
suis ni meilleur ni pire que n'importe qui. Si toutes ces choses sont ainsi,
pourquoi est-ce que j'imagine encore que tu vis, que tu es quelque part heureux
et en sécurité? Et pourrais-je vivre sans cette triste lâcheté?

La nuit dernière je rêvai que je reposais sur un champ, ou peut-être un parc. Je
posais ma tête sur ta jambes; c'est-à-dire sur te jambes et sur celles de
l'autre toi, celle qui était morte. Tes cheveux tombaient comme une cascade sur
mon front. Je voulais les embrasser, mais une paralysie total m'en empêchait. Je
suppose que ce rêve est, a sà manière, adèquat. Il existe deux toi. Mais je aime
seulement à une.

\pagebreak 

$\S$ Merci Seigneur, pour la musique que j'ai faite aujourd'hui avec V et M, et pour
avoir facilité ma situation de logement. J'ai bon espoir que Daniel vivra
dignement et je te demande, Jésus-Christ, de ne pas permettre que j'agisse mal.
Tu avais raison, Jésus-Christ, mon père, quand tu disais que je ne devrais pas
m'inquiéter pour demain. Je veux mettre ma vie entre tes mains. Tu es mon
refuge, car tu es la lumière de ma conscience.

\pagebreak 

$\S$ La nuit tombe profondément: Je sens le clair de lune avec\footnote{El
original es \textit{with}, no \textit{in}: el corazón es aquello \textit{con lo
que} se siente.} mon cœur. Encore une fois, je ne peux pas dormir. Ce qui me
dérange, ce sont ces perles, ces petits diables: les souvenirs. Bien que prenant
plusieurs formes, comme une eau mal intentionée, un visage, un nom, un homme
apparaît toujours. Ses yeux sont mille yeux, son visage est mille visages, son
âme est mille âmes. La violence et l'amour, la tranquillité et la vengeance,
l'obscurité et le silence, la fatigue et l'espoir: ce sont ses incarnations.

Nous avons beau essayer, aucun souvenir ne sombre entièrement dans l'oubli. Par
exemple, je me souviens maintenant d'eaux sombres sous une nuit d'été, de sable
gluant et désagréable sous les pieds. Maintenant, c'est la nuit: L et moi
nageons. Nous rentrons à la maison et nous découvrons qu'il n'y a plus
d'électricité. Nous faisons face à un dilemme familier: ouvrir les fenêtres pour
faire entrer l'air frais, mais avec les moustiques innombrables, ou conserver
les fenêtres fermées et résister à la chaleur. Mon père nous appelle depuis sa
chambre. Sur la petite terrasse, il a aligné trois matelas, les uns à côté des
autres. Nous sommes loin de la cité; aucune lumière artificielle ne trouble
l'air désert, et une myriade d'étoiles de toutes les couleurs embrase le ciel.
La brise est assez forte pour éloigner les moustiques. Alors que nous sommes
tous allongés sur la terrasse, sans couvertures ni oreillers, mon père commence
à parler: «Là, ce sont les Trois Maries; là-bas, c'est Orion...» Et ainsi nous
nous endormons.

~ 

~

\begin{itemize}
    \item «Ce» es pronombre demostrativo neutro: eso, aquello, lo...
    \item Pero «ces» es adjetivo demostrativo plural: estos, esos, aquellas, y
        necesita un sustantivo despues. «Estos pequeños diables» $\to$ «ces
        petits diables».
    \item Sujeto + avoir (conjugado) + beau + infinitivo.
    \item En «nous sommes tous allongés», «tous» significa «todos» o «todos
        nosotros». No modifica directamente al verbo, sino al sujeto «nous»:
        «todos nosotros estamos acostados». 
    \item «Là» significa «ahí», «allí», «ese punto». Se usa para señalar algo
        relativamente cercano o inmediatamente presente. «Là-bas» significa
        «allá», «más lejos», «en aquel lugar». En «Là, ce sont les Trois Maries
        ; là-bas, c'est Orion», padre señala primero una constelación cercana
        en el cielo visible y luego otra más distante.
\end{itemize}

\pagebreak 



\begin{forest}
[S
  [NP [Je]]
  [VP
    [V [me suis allongé]]
    [PP [sur le lit]]
  ]
]
\end{forest}

\ex
\begingl
\gla Je me suis allongé sur le lit //
\glb I REFL am laid.down on the bed //
\glft “I lay down on the bed.” //
\endgl
\xe
\noun{chien} \verbf{court} \prep{dans} le jardin 
\pagebreak  

La tombe dit à la rose :\newline
– Des pleurs dont l’aube t’arrose\newline
Que fais-tu, fleur des amours ? \newline
La rose dit à la tombe : \newline
– Que fais-tu de ce qui tombe\newline
Dans ton gouffre ouvert toujours ? \newline

La rose dit : – Tombeau sombre, \newline
De ces pleurs je fais dans l’ombre\newline
Un parfum d’ambre et de miel. \newline
La tombe dit : – Fleur plaintive, \newline
De chaque âme qui m’arrive\newline
Je fais un ange du ciel ! \newline


\begin{center}
\begin{forest}
[S
  [NP
    [Det [La]]
    [N [tombe]]
  ]
  [VP
    [V [dit]]
    [PP
      [P [à]]
      [NP
        [Det [la]]
        [N [rose]]
      ]
    ]
  ]
]
\end{forest}

\begin{forest}
[S
  [PP
    [NP
      [N [pleurs]]
      [CP
        [C [dont]]
        [S
          [NP
            [Det [l']]
            [N [aube]]
          ]
          [VP
            [NP [t']]
            [V [arrose]]
          ]
        ]
      ]
    ]
  ]
  [VP
    [V [fais]]
    [NP
      [Pro [tu]]
    ]
    [NP
      [N [fleur]]
      [PP
        [P [de]]
        [NP
          [Det [les]]
          [N [amours]]
        ]
      ]
    ]
  ]
]
\end{forest}

\end{center}

\pagebreak

Nous sommes tous, à toute heure, assiégés par des forces infinies qui nous
invitent à la quiétude, au conformisme et à l'idiotie. Dans tous les sujets sur
lesquels nous pensons, nous nous formons, par précipitation ou négligence, des
opinions sans mérite. Dans toutes les conversations, des paroles vides, des
paroles qui ne sont pas de notre cœur, veulent être dites. La main séductrice de
la paresse intellectuelle nous appelle sensuellement. C'est la lutte véritable.
Nous devons cultiver une intransigeance saine, une révolte adolescente. Nous
devons nous éloigner des conversations vaines. Nous devons échapper à la
frivolité. À chaque fois que nous sentons que nous perdons notre temps, nous
devons nous rappeler: dans ce paradis il y a des choses qui peuvent
m'émerveiller et que je dois découvrir. La vie est dans l'émerveillement.

La courbe de la vie commence à se dessiner sur le visage de mes contemporains.
Sur leurs visages je devine la trace abjecte ou adorable qu'ont imprimée leurs
désirs, leurs hontes, leurs péchés doux ou sinistres. Chez tant d'entre eux je
vois une complaisance déshonorante, un désintérêt obscène pour l'amélioration de
soi, un oubli délibéré que les étoiles existent. Ils se moquent de savoir que
dans le temps et l'espace s'embrasent des univers ou d'autres hommes, dont les
langues et les signes étant autres, donnent un sens singulier à l'aventure de la
vie. Ils commencent à boire en excès, à se prendre des libertés. Ils s'arrêtent
de lire. Ils s'arrêtent de penser. Quelques-uns s'arrêtent de sentir.

Dans chacun de mes cauchemars, il y a une teinte de vert pâle qui doit
ressembler à celle de ces vies qui errent d'une frivolité à l'autre, qui cachent
le fond fangeux de l'angoisse où le désir et l'amour se sont enterrés. Ils sont
comme les rivières dont les superficies turbulentes dissimulent un fond
immobile et marécageux. (Moi, dans ma retraite juvénile, je parais lent et
statuaire, mais je traverse des mondes millénaires et je sème sur mon fond boueux
de belles algues et des serpents).

\pagebreak 


Ce n'est pas souvent considéré, mais un des points critiques de la théorie
darwinienne est l'insubstantialité morale. Pour l'éthos chrétien, du moins
d'après l'interprétation traditionnelle, le bien et le mal sont des substances
actives. (La doctrine de la \textit{privatio boni}, indépendamment de ce que nous
pensons, est artificielle et tardive.) Mais, d'un point de vue métaphysique, le
monde darwinien est un monde insubstantiel, et notre manière de considérer
certaines choses comme bonnes ou mauvaises est un phénomène biologique que la
neuroscience moderne réduit, correctement à mon avis, au cerveau.

Je ne veux pas discuter de ces choses; je voudrais plutôt montrer comment Redon a
exprimé ces problèmes dans son œuvre. La raison est triple.

Premièrement, Redon étudia la théorie de l'évolution avec un certain degré de
détail. C'était un ami d'Armand Clavaud, un botaniste qui l'initia aux lectures
de Darwin et Lamarck, mais aussi à celles de Baudelaire, Poe, et certains poètes
hindous. Redon se fascina pour ce que l'on pourrait appeler le problème de
l'origine, et produisit un album magnifique nommé précisément \textit{Les
Origines}. En 1881, il vit à Paris une exposition des aborigènes captifs de
la Terre de Feu et dit:

\begin{quote}
They provided me with a dream of primitive life, nostalgia for the pure and
simple existence of our beginnings. I have never felt with as much force the gap
our own nature spans between the crawling beast and our higher destiny. . .
.What a poem, this quite perfect organism, emerging from the primeval slime to
stand beside us and stammer out the first stanzas of a universal hymn.
\end{quote}

\begin{figure}[htbp]
    \centering
    \includegraphics[width=\textwidth]{../Images/LesOrigens.png}
    \caption{Les Origines}
    \label{fig:lesorigens}
\end{figure}

Cette citation révèle sa fascination, bien que mélangée aux préjugés propres
de son temps, pour la «vase primitive» et l'histoire évolutive de l'homme.

Deuxièmement, Redon comprenait les problèmes posés par la théorie de l'évolution
de manière plus sophistiquée que ses contemporains. L'Européen moyen
s'inquiétait du passage d'une forme à une autre, un problème réel mais mineur
d'un point de vue scientifique et spirituel. La mesure dans laquelle l'Europe se
préoccupait de cette question se révèle dans les nombreuses caricatures
représentant des singes qui se transforment en hommes, satires qui au fil du
temps se transformèrent en satires d'elles-mêmes. Mais Redon était fasciné par
un problème plus fondamental et complexe; à savoir, le passage de l'informe à ce
qui a une forme. C'est le problème du xάος hésiodique, que Darwin ressuscita
et qui est plus difficile à résoudre et à caricaturer. Par exemple, dans \textit{Les
Origines} nous voyons l'œuvre \textit{Quand s'éveillait la vie au fond de la
matière obscure}, dans laquelle une créature est en train de surgir précisément d'une
vase primitive, tandis que, sur le fond obscur et depuis l'espace, des têtes
flottantes s'approchent de la Terre. D'après Sven Sandström, le premier expert à
organiser \textit{Les Origines}, la lecture de cette œuvre est temporelle: «When
the seed from space arrives on earth life starts in early states to develop from
matter."


\begin{figure}[H]
    \centering
    \includegraphics[width=\textwidth]{../Images/Redon2.jpg}
    \caption{Quand s'éveillait la vie au fond de la matière obscure}
    \label{fig:redon2}
\end{figure}

Il y a un autre example dans l'œuvre \textit{ Le Polype difforme flottait sur
les rivages, sorte de cyclope souriant et hideux}, aussi de \textit{Les
Origines}, dans laquelle nous voyons une «polype» (le mot fait référence à
Polyphème). La créature n'est ni animal ni homme. Ses traits grotesques évoquent
le primitif et ce qui n'a pas encore de forme, mais qui est en train de
l'acquérir. Mais, comme une étude a indiqué [1], le «polype» est aussi 


\begin{quote}


«... un organisme unicellulaire qui faisait l'objet de nombreuses investigations
scientifiques à l'époque, non seulement en tant que l'une des premières formes
de vie, mais aussi parce qu'il était étrangement inclassable. Ni tout à
fait plante, ni tout à fait animal, il fut désigné par les naturalistes comme un
emblème de l'informe et de l'indéterminisme»

\end{quote}



\begin{figure}[H]
    \centering
    \includegraphics[width=\textwidth]{../Images/Redon4.png}
    \caption{Le Polype difforme flottait sur les rivages, sorte de cyclope souriant et hideux}
    \label{fig:redon2}
\end{figure}

Finalement, malgré sa compréhension de la théorie de l'évolution et son désir
d'imaginer le passage de l'informe à la complexité de la vie comme quelque chose
de matériel, Redon pressentait une force spirituelle sous-jacente à la vie. Il a
même dit, dans ses \textit{Confidences d'Artiste}, qu'Armand Clavaud était un
admirateur fervent de Spinoza, et il n'est pas difficile d'imaginer que Redon le
lisait aussi. Il dessina un Christ, mais appela l'œuvre \textit{Buda}, en
suggérant que toutes les religions émanent d'une source commune. Bien qu'immergé
dans le cercle symboliste ---il était habitué des réunions de Mallarmé---, son
intérêt scientifique lui fit comprendre comme aucun autre les principes
spinozistes qui sous-tendent la théorie du double aspect. Tandis que la plupart
des symbolistes pressentaient dans la réalité sensible un secret esthétique soit
merveilleux, soit décadent, Redon voyait un arcane spirituel qui non seulement
ne contredisait pas la science, mais la sublimait.

\begin{figure}[H]
    \centering
    \includegraphics[width=\textwidth]{../Images/Redon3.jpeg}
    \caption{Buda}
    \label{fig:redon2}
\end{figure}

La vision de la science de Redon était \textit{correcte} au sens usuel du terme:
à savoir, elle est en accord avec mes jugements et préjugés. Blague à part, la
vision spinoziste du monde est, il me semble, la seule possible face à tout ce
que nous avons appris de la neuroscience. Comme c'est souvent le cas, l'art
surpasse la science pour exprimer la nature des choses. La plupart des
références à Redon en ligne le décrivent comme un fou, comme un homme qui
coloria des visions étranges et inexplicables. Mais il ne fit rien de plus que
se poser les bonnes questions face aux découvertes de son temps et les sublimer
avec un art exquis.

Je voudrais commenter un dernier aspect de l'art de Redon. Jusqu'ici, j'ai parlé
de Redon comme quelqu'un qui comprit certains problèmes propres à son époque et les
exprima avec une imagination et une sensibilité particulières. Mais je n'ai rien
dit à propos du problème de l'insubstantialité morale : j'ai seulement raisonné
sur le problème de l'origine, et j'ai suggéré vaguement une ontologie spinoziste
comme solution à l'apparent vide ouvert par Darwin et la science en général
[[2]]. Quelle que soit la réponse que Redon trouva au problème moral, nous
l'ignorons; mais dans une de ses dernières œuvres je comprends qu'il en a trouvé
une. Car le plus grand signe de richesse morale et spirituelle est la capacité
de comprendre le mal avec compassion et tendresse, et c'est précisément ce que
Redon montre dans l'huile \textit{Le Cyclope}.

Dans \textit{Le Cyclope}, Redon explore le motif classique du cyclope Polyphème
et son amour pour Galatée. Bien qu'Ovide, dans le livre XIII des
\textit{Métamorphoses}, ait su décrire «cette créature impitoyable, terrifiante
même pour les forêts, que nul étranger n'a jamais croisée impunément», avec
une certaine tendresse, la plupart des représentations de Polyphème dans l'art
sont bestiales et impies. Elles représentent le monstre, la brute, l'idiot.
Redon est l'un des rares qui devina dans l'amour de Polyphème une force si
tendre, si pure, si innocente, qui constituait peut-être une qualité rédemptrice.
Dans \textit{Le Cyclope} nous voyons Polyphème en train de garder Galatée alors
qu'elle dort. Son visage, bien que bizarre, est aimable et tranquille. Il la
voit comme quelqu'un qui surveille le sommeil d'un être cher. Et tout arrive dans une
scène champêtre et pacifique, où les couleurs ---choisies dans le style typique du
Redon tardif--- transmettent un bonheur paisible.

C'est suffisant. J'espère avoir montré deux choses. D'un côté, que Redon n'était pas
un fou qui coloriait des fantaisies inexplicables, mais un artiste profondément
intéressé par l'idée d'exprimer certains problèmes posés par le darwinisme et la science
de son temps. De l'autre, que son intérêt pour la science non seulement ne l'a
empêché de poursuivre avec une forme chrétienne et singulière de symbolisme, mais
que ce mélange dotait son art d'une originalité et d'une pitié sublimes.




~


[1]: https://www.researchgate.net/publication/345364261_Into_the_Primeval_Slime_Body_and_Self_in_Redon's_Evolutionary_Universe

[2]: Je pars du principe que le matérialisme ontologique est en totale
contradiction avec la conception scientifique du monde, une thése que je n'ai
pas le temps de développer.

\pagebreak

$\S$ En ce moment, je suis à la maison de campagne. Nous sommes le 25 mai,
et au rez-de-chaussée les gens continuent de célébrer l'indépendance de
l'Argentine. Moi, d'autre part, je suis dans ma chambre : fatigué, somnolent,
désireux de faire quelque chose d'enrichissant avec mon temps. Quelle misère de ne
pas pouvoir jouir de la compagnie du prochain!

Une certaine lascivité brûle en moi. Un certain éloignement, une suspension
du ton affectif, une quiétude sans paix. Au loin la lagune, terriblement
immobile, reflète les nuages du ciel froid. Le chant d'un hurleur noir éclipse
les voix humaines qui seraient autrement un acte de violence contre le mutisme du
monde. La rumeur de ce que l'on pourrait appeler le crépuscule, bien qu'il ne
soit ni rouge ni doré, me torture en me rappelant que la lumière du Christ brûle
péniblement en moi, à peine perceptible, à peine vivante, comme une luciole sur
le point de devenir une ombre.

Je veux retourner mon cœur pour que le sable tombe à nouveau, car en ce moment il
ne me reste plus de temps.

\pagebreak 

Je suis l'eau ténébreuse de la rivière,\newline
dont la surface est ennemie des lucioles,\newline
en laissant une cendre sans lumière\newline
au cours de cette vie où tu t'étioles.

Ô, retourner mon cœur, pour que le sable \newline
tombe à nouveau et les nuits où je te vis \newline
soient immarcescibles, impérissables, \newline
comme ta mort est longue et infinie...!

\pagebreak 

$\S$ Dans la vie, il y a deux types de personnes douteuses: ceux qui n'ont pas
de principes et ceux qui en on trop. La laxisme et l'inflexibilité sont
également désagréables.

Au cours de ma vie, j'ai embrassé huit principes. Je voudrais les exprimer
brièvement, puisque le temps qu'il me reste est court. L'antécédent familial
pronostique une mort jeune pour moi, et, du moins selon lui, j'ai traversé
probablement plus de la moitié de ma vie. Mais même sans cette donnée médicale,
j'avoue qu'un pressentiment de mort spirituel et dérangeant vit en moi. Je fais
des rêves où sa main dangereuse me caresse. Je vais bientôt mourir, et, avant
que mon heure vienne, je voudrais tracer l'éphémère mémoire du monde avec les
huit étoiles qui me guident. Ce sont:


\small
\begin{quote}

$(1)$ Non-résistance au mal.

$(2)$ Opposition à toute autorité terrestre et à tout pouvoir concentré,
considérés comme intrinsèquement mauvais.

$(3)$ L'austérité volontaire, ou plutôt l'abstention volontaire de la poursuite
de la richesse, du pouvoir ou du luxe.

$(4)$ S’abstenir de juger les autres, comme l’enseigne Matthieu 7:1-5, et de
toute arrogance et vanité.

$(5)$ Tout pardonner, comme l'enseigne Matthieu 18:21-22.

$(6)$ Se ranger du côté de ceux qui ont faim, qui sont démunis, qui sont
emprisonnés, qui sont soumis et qui sont faibles, et aller à leur rencontre.

$(7)$ Ne jamais, sous aucun prétexte, se permettre de limiter la liberté d'un
autre homme ou d'une autre femme, que ce soit sur le plan économique, politique,
religieux, sexuel, etc.

$(8)$ Ne jamais s'inquiéter : ni de sa propre vie ni de ce que demain pourrait
apporter, comme l'enseigne Matthieu 6:25-34.

\end{quote}
\normalsize

\pagebreak 

$\S$ Le temps qu'il me reste est court. L'antécédent familial pronostique une
mort jeune pour moi, et, du moins selon lui, j'ai traversé probablement plus de
la moitié de ma vie. Mais même sans cette donnée médicale, j'avoue qu'un
pressentiment de mort spirituel et dérangeant vit en moi. Je fais des rêves où
sa main dangereuse me caresse. Je vais bientôt mourir. Je suspecte que
Jesus-Christ et moi avons fait un pact.

~

Quand j'avais vingt ans, ma destinée m'avait éloigné de ceux que j'aimais, dans
une ville inconnue, en conflit avec tous. Mon amour intime mourut: loin de moi,
elle s'est donné la mort et je ne pus ni voir son corps ni lui dire adieu. Nous
avions été chéris ou meilleurs amis par intermittence depuis l'âge de douze
ans... Ma famille était défaite, mes études, interrompré pour le pandemie, ne
signifiaient rien pour moi. J'étais seul parfaitement. 

Ainsi, immergé dans la nuit de la vie, je décidai que mon existence devait finir
et je voulus la mort. Mais, alors que j'arrivais à l'heure convenue, je vis avec
un œil intérieur la figure d'un homme, entouré de lumière, qui conduisait un
agneau sur une cime derrière laquelle le soleil déclinait. Et lorsque je vis
cette image, la paix du monde m'inonda. Instantanément, mon désir de mourir
s'est dissipé et ma dépression, qui dura une année entière, donna graduellement
lieu à l'espoir.

Aujourd'hui, je suis heureux, du moins la plupart du temps. Mais je m'accroche à
la lumière de mon Christ péniblement quand je pense à cet amour qui n'existe
plus.


\pagebreak 

$\S$ J’ai rencontré aujourd’hui une personne avec une histoire similaire à la
mienne. Elle avait trente-tois ans. Elle me semblait vulnérable, bien qu’elle cachât sa
vulnérabilité sous une surface dure et catégorique. Comme toutes les personnes,
elle avait besoin d’amour, de tendresse, de loyauté et de gentillesse. Beaucoup
de souffrance dans sa vie. Elle traversait une crise d’identité non négligeable.
Une séparation la laissait à la recherche de son propre sens. Elle avait des
croyances bizarres : l’astrologie, les énergies, ce genre de spiritualité qui
exige très peu d’effort. Je crois qu’elle est généreuse, et peut-être brisée.

\pagebreak 

Dans la \textit{Lettre IX} de Voltaire, \textit{Sur le Gouvernement}, j'ai
appris certain choses de la langue française que je ne connaissais pas. 

$(1)$ «Ils s'en vantent», dit Voltaire des Anglais, mais on typiquement dit «Ils
se vantent \textit{de} qualque chose». Quand un verb transitif se construit avec
une préposition, «$v$ de $o$», on peut aussi construire la phrase «en $v$», où
«en» remplace «de $o$» avec un objet direct antérieur. La phrase complète de
Voltaire est: 

\begin{quote}
Il est vrai qu'avant et après Guillaume le Conquérant les Anglais ont eu des
parlements; ils s'en vantent, [...]
\end{quote}

\begin{quote}
    Es cierto que antes y después de Guillermo el Conquistador los ingleses
    tuvieron parlamentos; se jactan \textit{de ello}, [...]
\end{quote}

o \textit{de lo cual se jactan}, etc.

$(2)$ «Les papes se mîrent a la tête», il dit, une phrase qui signifie 
«los se pusieron a la cabeza», dans le sense de «se pusieron a dirigir», «se
pusieron a liderar», etc. 

$(3)$ «Les rois alors n'étaient point despotiques», dit Voltaire. La
construction «ne $v$ point» est une forme de négation qui introduit une nuance
de négation plus forte, similaire a «no $v$ en absoluto».

$(4)$ Il continue: «mais les peuples n'en gémissaient que plus dans une servitude
misérable». La construction «$s$ ne $v$ que ...» exprime que le sujet de la phrase
ne fait plus que $v$. En español, se podría traducir como «no más que $v$», «$s$
no hace sino $v$». 

«$S$ ne dit que le nomme de $J$»: «$S$ no dice más que el nombre de $J$», «$S$
no hace sino decir el nombre de $J$», etc.

Voltaire use cette construction avec «plus», ce qui exprime que les peuples
gémissaient plus que jamais, ou que les peuples gémissaient plus que dans
d'autres circonstances, etc. En español, se podría traducir como «pero los
pueblos no gemían sino más en una servidumbre miserable», «pero los pueblos no
hacían sino gemir más en una servidumbre miserable», etc.

\pagebreak 

$\S$ Dans l'essai \textit{Sur le Gouvernement}, Voltaire, inspiré par son séjour
dans ce pays, décrit l'histoire politique de l'Angleterre et la compare avec
celle de la France. Parce que certains éléments de ce discours sont révélateurs
de sa conception de l'État et les problèmes de son temps, leur analyse est
profitable.

Les critiques de Voltaire se concentrent principalement sur l'Église et la
noblesse. Par exemple, à ses yeux le pouvoir des évêques, moines et papes est
une continuation du pouvoir ancien des druides, figures qui sont assez
primitives et païennes. Cela peut paraître peu, mais à une époque où le pouvoir
de l'Église était si superlatif, la concevoir ainsi n'est pas rien. L'Église est
vue comme une force politique motivée par l'accumulation de l'argent et non
comme l'incarnation institutionnelle de la volonté divine. Et la noblesse,
d'autre part, est vue comme une force émergée des sauvages qui se
firent monarques et qui «disputaient souvent avec leur roi les dépouilles des
peuples». Ce n'est pas une vision très flatteuse, c'est le moins qu'on puisse
dire.

En général, Voltaire est plus critique envers l'Église qu'envers l'État.
Voltaire méprise toute condition dans laquelle un peuple est subjugué par de
multiples maîtres, mais considére le gouvernement d'un roi unique, limité par un
régime constitutionnel, comme un bien. En décrivant les barons, les évêques, et
les papes comme des vampires et des tyrans, il demande:

\begin{quote}
Et n'est-ce pas un
bonheur pour le genre humain que l'autorité de ces petits brigands ait été éteinte
en France par la puissance légitime de nos rois, et en Angleterre par la puissance
légitime des rois et du peuple?
\end{quote}

Le «concert entre les Communes, les Lords et le Roi» est vu comme la solution du
problème de la tyrannie «goth et vandale» de laquelle le pouvoir de l'Église et
la noblesse émergent.

La vision de Voltaire, cependant, n'est pas naïve. Cela est clair surtout dans
sa description de la \textit{Grand Charte}. Il voit ce document, largement
considéré comme le début des libertés Anglaises, avec certain scepticisme. Considère,
par exemple, le passage suivant:

\begin{quote}
    Les barons forcèrent Jean sans Terre et Henri III à accorder cette fameuse
    charte, dont le principal but était, à la vérité, de mettre les rois dans la
    dépendance des lords, mais dans laquelle le reste de la nation fut un peu
    favorisé, afin que, dans l'occasion, elle se rangeât du parti de ses
    prétendus protecteurs.
\end{quote}

Cette petite description réaliste et pragmatique de la charte peut très bien être
une description de tous les lents progrès de l'humanité.

De plus, on devrait comprendre précisément ce que Voltaire admirait dans le
gouvernement des Anglais. Il souligne que les seigneurs et les évêques n'ont pas
le pouvoir de faire ni de modifier les lois; que chacun donne non selon sa qualité
mais selon son revenu, et qu'il y a une taxe réelle sur les terres.

Finalement, ce que Voltaire dit à propos des causes réelles de la liberté en
Angleterre n'est pas de moindre importance : c'est dans les fissures à la
surface de la tyrannie que la liberté naît, comme la mousse dans les rochers :

\begin{quote}
Heureusement, dans les secousses que les querelles des rois et des grands
donnaient aux empires, les fers des nations se sont plus ou moins relâchés; la
liberté est née en Angleterre des querelles des tyrans.
\end{quote}

De nouveau, ce n'est une observation qui s'applique seulement à l'Angleterre. 

\pagebreak 

J'écris dans Philadelphia, cet endroit étrange où on voit collègiens heureux
avec toutes les opportunités imaginables, et, à côté, des spectres consumé pour
la pauvreté, les opïodes, et la folie. La vie américaine, en supposant qu'on est
o peut être aveugle, est facile; facile, je veux dire, pour ceux qui ont encore
un corp et un esprit, et qui ne sont pas devenus de fantômes irréversibles.

En traversant l'université, qui est très belle, j'ai la même sensation que
j'avais autrefois : les étudiants n'ont pas la moindre idée de combien ils sont
privilégiés. J'en parle avec ma mentor, car elle pense la même chose : tout est
considéré comme acquis. Même celles qui disent savoir, et celles qui désirent
savoir, ne peuvent pas le faire. Je pense à mes camarades qui vont à notre
université avec des chaussures déchirées et qui ne peuvent pas se permettre de
sauter un repas au restaurant universitaire, lequel coûte l'équivalent de 1,3
dollar ; ils ont beaucoup de capacités et de potentiel, et notre gouvernement
les méprise toujours plus. Qu'est-ce qu'ils font pour mériter cela ? Rien. Et
qu'est-ce que les Américains font pour mériter leur fortune ? Rien.

Je ne méprise pas les Américains : je célèbre les opportunités qu'ils reçoivent
ici, parce que je crois que c'est un droit qu'un homme puisse étudier. C'est un
droit parce que la liberté l'est, et l'éducation est une condition nécessaire
pour la liberté, autant que l'ignorance l'est pour l'esclavage. Je désirerais
seulement que tous les hommes aient l'éducation scientifique et culturelle que
je vois ici ; je reconnais que cette inégalité est arbitraire, et je m'indigne
devant l'injustice de voir des jeunes plus talentueux que moi, et que beaucoup
de ceux que je vois ici, limités par des conditions qui, en tout cas, les
précèdent de cinq cents ans.

Par ailleurs, les Américains me paraissent être des personnes très distraites,
une caractéristique que, dans une certaine mesure, ils partagent avec les
Argentins. On entend peu de conversations sur des sujets politiques,
philosophiques ou culturels, y compris dans l'université. Les conversations
intellectuelles traitent en général de sujets très spécifiques à la discipline
de chacun, et l'introduction de propositions de nature fondamentale ou
philosophique est vue avec un certain étonnement. En général, leur système
éducatif aspire à former des spécialistes, et non des penseurs, un modèle qu'ils
reconnaissent explicitement et qui sans doute a ses avantages. Et, bien que les
spécialistes ici soient extraordinaires, cette conception de ce qui signifie
être éduqué se conjugue avec un certain manque d'esprit critique, et avec une
culture qui valorise la consommation, le pouvoir et l'argent avant tout. Alors,
nous avons précisément ce que je vois : des personnes qui, même en tant
qu'éduquées ou intelligentes, peuvent vivre sans jamais se poser les questions
les plus urgentes de leur société, sans jamais admirer la beauté d'un poème, et
qui, surtout, peuvent voir un homme dormir tremblant de froid au pied d'un
gratte-ciel et ne pas vouloir mettre le feu à quelque chose.

Mais je ne veux pas être injuste. Les Américains de Philadelphia ont beaucoup de
qualités admirables. Surtout, j'admire leur tolérance et la diversité qu'elle
engendre. Dans cette ville, on entend des centaines de langues différentes et on
voit des personnes de toutes les ethnies imaginables ; là, le chrétien côtoie le
musulman ; là-bas, le juif et l'athée sont voisins, et personne ne se sent
menacé par la différence de l'autre. L'Amérique primitive, celle qu'on associe à
la bannière confédérée, n'existe pas ici, ou du moins elle n'est pas visible.

Il faut dire, comme toujours, que les apparences sont trompeuses. Il y a un
revers à ma description, et c'est que Philadelphie est une ville très ségréguée.
C'est en raison du « redlining », une injustice que le gouvernement a
institutionnalisée dans les années 30 et qui a laissé des traces que nous voyons
encore aujourd'hui. La criminalité se concentre dans ces zones, curieusement là
où la pauvreté se concentre aussi... Ces citoyens sont privés de l'opportunité
de profiter des vastes ressources de la ville, indépendamment du fait que leurs
concitoyens les traitent avec respect et sans discrimination explicite. La
coexistence de ce problème structurel avec la tolérance dans les interactions
quotidiennes est un paradoxe seulement apparent, car il reflète précisément ce
que je disais plus tôt : même les personnes les plus éduquées, et par conséquent
les plus tolérantes, peuvent vivre en paix avec les injustices de leur société,
parce que leur modèle de vertu n'a rien à voir avec la solidarité sociétale, ni
avec aucune des vertus typiquement associées au christianisme. Dans cet endroit
fébrile, les hommes ne se voient pas les uns les autres comme des frères ; leur
dieu est l'argent, et leur péché le plus fatal est de ne pas se savoir ignorants
et de vivre dans l'illusion d'une liberté irréelle.

Mais les Américains ne sont ni idiots ni soumis : s'ils savaient que leur
liberté est vaine, que bien qu'ils aient des centaines de bêtises à choisir,
dans ce qui compte ils n'ont guère plus d'un ou deux choix, leur gouvernement ne
pourrait pas subsister sans une forme totale, voire insoutenable, de répression.
C'est pourquoi, comme dans aucun autre lieu au monde, on voit ici un appareil si
sophistiqué de propagande et de distraction systématisées. Tout le monde parle
de choses : d'un électroménager, d'une voiture, des prix, de superficialités. Ce
peuple n'est pas ainsi par nature : leur gouvernement, et j'inclus dans cette
catégorie le capital concentré, les prive de développer les qualités humaines
qu'ils auraient autrement. La preuve de cela est simple : si les Américains
étaient par nature enclins à la consommation matérielle, il ne serait pas
nécessaire de dépenser autant d'argent pour les faire consommer, et la
publicité, industrie qui n'est pas sans importance ici, n'existerait pas.

Le contrat dans ce pays est le contrat oligarchique typique: una vie
confortable en échange d'une certaine passivité. Comme l'a dit James Madison,
\textit{founding father} et le «père de la constitution», l'impératif est ici de
«to protect the minority of the opulent against the majority». Comme on le sait,
il y a deux stratégies pour atteindre ce but: le bâton, ou cette combinaison de
confort superficiel et propagande. Dans l'Amérique latine, nous connaissons la
première, même si nous sommes en transition vers la seconde; aux États-Unis,
le bâton existe aussi, mais il est utilisé plus rarement.

On pourrait penser que j’exagère, que précisément parce que la vie américaine
est relativement facile, il est justifié qu’ils ne vivent inquiets pour rien.
Mais rien n’est plus éloigné de la vérité. En interne, les problèmes sont
visibles. Considérons seulement le sans-abrisme à Philadelphie par rapport à
celui de Córdoba, en Argentine : le nombre d’habitants est plus ou moins le
même, mais la première a dix fois plus de sans-abri que la seconde, selon les
rapports officiels. Mais lorsqu’on considère la politique étrangère, la
situation est encore pire : les Américains savent, car il est impossible de
l’ignorer, que leurs impôts sont en grande partie dirigés vers la guerre, la
destruction, l’invasion et la mort. J’imagine qu’eux-mêmes ne croient plus à la
fable de l’exportation de la démocratie. Mais je n’ai entendu pas un seul mot
sur l’Iran, par exemple. Pour eux, avoir agressé un pays, avoir causé la mort de
près de 170 enfants avec un missile à Minab, avoir frôlé la guerre nucléaire,
sont des faits devenus banals. Dans toute société saine, de tels événements
provoqueraient au moins des manifestations de masse. Ici, soit on les ignore,
soit on préfère ne pas en parler.

\pagebreak 

$\S$ Quand je me sens triste sans raison, je sais précisément quelle est la
sensation qui m'attriste: ma vie ---présente, passée et future--- se présente à
ma conscience comme une œuvre d'Edward Hopper. Cela est tellement vrai, que
quand je découvert cet artiste, ma réaction immédiate fut de dire : «enfin ma
tristesse a trouvé son visage».

Je suppose que l'unique forme de tristesse que j'ai connu est la solitude. Le
gens disent qu'il faut jouir de sa propre solitude, que l'homme fort a paix en
elle. Si tel est le cas, je ne suis pas fort. Dans mon moment plus obscure,
Travis Bickle fut le personnage qui exprimé au mieux ma vie intérieur.

\pagebreak 

$\S$ J'ai un rêve récurrent si curieux et bête qu'il me semble amusant de le
partager. Je fais ce rêve au moins trois fois par an. Bien qu'il présente
certaines variations, le fond est toujours le suivant : Je traverse une boutique
de pipes à tabac et j'admire leurs designs exquis. À l'occasion, je découvre
plusieurs pipes avec le design appelé \textit{Jade}. Ces pipes sont reluisantes,
et leur couleur est le vert profond du mineral qui leurs donne son nom. (Mon
père m'acheta l'une de ces pipes quand j'eus seize années, et peut-être ai-je
ressenti ce moment-là comme un rite de passage). Le rêve n'est rien de plus que
cela: la contemplation de ces pipes, vertes et admirables, aux formes que le
rêve crée et que le monde n'a jamais vues. 







\pagebreak 

À vrai dire, la table d'Ivan, près de la fenêtre, était simplement protégée par
un paravent contre les regards indiscrets. Elle se trouvait à côté du comptoir,
dans la première salle, où les garcons circulaient à tout moment. Seul un vieux
militaire en retraite prenait le thé dans un coin. 
Dans les autres alles, on
entendait le brouhaha habituel à ces établissements: des appels, les bouteilles
qu'on débouchait, le choc des billes sur le billard.
Un orgue jouait. Aliosha
savait que son frère n'aimait guère les cabarets et n'y allait presque jamais.
Sa présence ne s'expliquait donc que par le rendez-vous assigné à Dmitri.

\begin{quote}
    Observaciones: 

    \begin{itemize}
        \item  Ne $v$ donc que ... puede usarse en vez de Ne $v$ que.
        \item Acordate que «on» como sujeto puede ser «uno», «they» o «la
            gente», «alguien», o incluso «nosotros» en francés moderno. En este
            caso, sería un «they» muy impersonal, y en castellano diríamos «las
            botellas que se descorchaban».
    \end{itemize}

\end{quote}

--- Souvent je me suis demandé s'il y avait au monde un désespoir capable de vaincre
en moi ce furieux appétit de vivre, inconvenant peut-être, et je pense quíl nén
existe pas, avant mes trente ans, tout au moins. Cette soif de vivre, certains
moralistes morveux et poitrinaires la traitent de vile, surtout les poètes. Il
est vrai que c'est un trait caractéristique des Karamazov, cette soif de vivre à
tout prix ; elle se retrouve en toi, mais pourquoi serait-elle vile ? On veut
vivre, et je vis, même en dépit de la logique. Je ne crois pas à l'ordre
universele, soit ; mais j'aime les tendres pousses au printemps, le ciel bleu,
j'aime certaines gens, sans savoir pourquoi. J'aime l'heroïsme, auquel j'ai
peut-être cessé de croire depuis longtemps, mais que je vénère par habitude.
[...] Comprends-tu quelque chose à mon galimatias, Aliocha ? conclut-il dans un éclat de rire.

— Je comprends trop, Ivan ; on voudrait aimer par le cœur et par le ventre, tu
l’as fort bien dit. Je suis ravi de ton ardeur à vivre. Je pense qu’on doit
aimer la vie par-dessus tout.

— Aimer la vie, plutôt que le sens de la vie ?

— Certainement. 
L’aimer avant de raisonner, sans logique, comme tu dis ; alors seulement on en comprendra le sens. Voilà ce que j’entrevois depuis longtemps. La moitié de ta tâche est accomplie et acquise, Ivan : tu aimes la vie. Occupe-toi de la seconde partie, là est le salut.

— Tu es bien pressé de me sauver ; peut-être ne suis-je pas encore perdu. En
quoi consiste-t-elle, cette seconde partie ?

— À ressusciter tes morts, qui sont peut-être encore vivants. 

...

Souvent je me suis demandé s’il y avait au monde un désespoir capable de vaincre
en moi ce furieux appétit de vivre, inconvenant peut-être, et je pense qu’il
n’en existe pas, avant mes trente ans, tout au moins. Cette soif de vivre,
certains moralistes morveux et poitrinaires la traitent de vile, surtout les
poètes. Il est vrai que c’est un trait caractéristique des Karamazov, cette soif
de vivre à tout prix ; elle se retrouve en toi, mais pourquoi serait-elle vile ?
Il y a encore beaucoup de force centripète sur notre planète, Aliocha. On veut
vivre, et je vis, même en dépit de la logique. Je ne crois pas à l’ordre
universel, soit ; mais j’aime les tendres pousses au printemps, le ciel bleu,
j’aime certaines gens, sans savoir pourquoi. J’aime l’héroïsme, auquel j’ai
peut-être cessé de croire depuis longtemps, mais que je vénère par habitude


\pagebreak 
$\S$ Je pense alors qu'il est possible que je n'aie pas le sense de l'humour.
Oui, en compagnie des autres je suis relativement charismatique, d'accord avec
certains, même heureux. Mais, comme le dit un certain poème que j'écris depuis
longtemps, et s'ils me voyaient quand je suis seul ? 

\begin{quote}
    \textit{Et s'ils me voyaient dans ces nuits désertes,}\newline
    \textit{quand, sans sommeil et en cherchant quelque chose de
    précieux},\newline
    \textit{je doute des vérites sur lesquelles ma foi fut construite}?
\end{quote}

Je me découvre étrange : un scientifique confus qui trouve Dieu, qui se demande
en permanence de si son cœur n'est donc que ténèbres, qui se dit : 
«tu dois aimer la vie», et parfois le fait. 

Dostoievski, du moins à travers de la voix de Alyosha, dit qui on dois aimer la
vie plutôt que le sens de la vie,

\begin{quote}
 L’aimer avant de raisonner, sans logique, comme tu dis ; alors seulement on en comprendra le sens. 
\end{quote}

C'est il que je veux faire : commencer à l'amour pour arriver à le sens. C'est
possible ? 

Le «Grand Inquisitor» est persuasive. Le soif de vivre est universele, le soif
brute, inevitable ; mais est plus concevible que l'homme n'eau pas capable
d'aimer la vie, même son próchême. Liberté : pourquoi ? Mais Alyosha comprendre
correctement que le discourse d'Inquisitor n'est donc que le plus grand
(\textit{elogio}) possible. 

\pagebreak 

$\S$ Il est une chose curieuse : avant, je pensais que le meilleur de moi
fleurissait quand j'étais en compagnie des autres. Seul, je me découvrais
mélancolique, morne et taciturne. Mais, alors, je comprends que précisément
lorsque je suis seul, et précisément depuis cette bataille contre ma tristesse,
ma vie dégage sa lumière la plus salutaire et secrète. 

La bataille, il faut bien le dire, est cruelle, et peut-être qu'il y a un
certain profit dans la considération de sa nature. Mais, il n'est pas
clair par où on doit commencer : peut-être, je peux commencer ainsi : Il n'y a
rien de plus étrange pour moi que ceux qui sont à la fois heureux et simples.
Même quand je suis heureux, cet état n'est qu'un équilibre fragile, comme
celui d'une boule sur le maximum local d'une courbe concave, au point où la
dérivée est exactement zéro. Les forces du monde, extérieures et intérieures,
l'incitent à trembler, bien que dans un rayon limité, violemment. Et il est
impossible de nier que ce genre de personnes existe : des gens de bien qui
ne sont pas assiégés par le doute, l'agonie, l'affrontement entre le bien et le
mal. J'ai pensé parfois, en les regardant : « Le pauvre ! Qu'est-ce que tu feras
quand la vie ouvre sa gueule insensible devant ton innocence ? » Mais c'est une
question idiote : qu'est-ce que je fais ? En outre, cette classe de personnes, il
faut le dire, est belle, même terriblement. Ce qui me perturbe est simplement
que, face à leurs visages, je découvre que mon propre visage est une toile
décatie sur laquelle la vie a gravé un dessein jamais désiré par moi. 
Notez que je ne dis pas, et ne prétends jamais dire, que je ne suis pas heureux :
plutôt, je veux dire simplement que le bonheur est pour moi comme l'eau coulante
et éphémère d'un ruisseau, plutôt que l'impassible surface d'un lac. 

C'est alors ce que je veux dire : 
les personnes qui sont à la fois simples et heureuses me font peur. Devant leur hédonisme ---elles ne peuvent être
qu'hédonistes--- leur courage, leur insouciance, je vois le pire en moi
et je sens comment mon rêve idiot d'être compris s'écroule.

Comme je peux être heureux avec une âme déchirée ; non pas misérable, mais
déchirée ! Mais face à un de ces cerfs immaculés et vierges, quel gouffre,
quel reflet infernel, quelles ténèbres !















\end{document}



